\chapter{Introduction}
\label{ch:intro}

Aircraft turnaround is the period between an aircraft arriving at its parking position and leaving for its next flight. Ground teams unload baggage, clean and refuel the aircraft, and prepare it for boarding. Some activities proceed in parallel, whereas others require preceding work to finish. Flights also compete for shared stands, cleaning teams, fueling vehicles and pushback tugs. A delay in one service can therefore affect both its own turnaround and the resources available to other aircraft.

This thesis studies those interactions using an object-centric simulation informed by available Berlin Brandenburg Airport (BER) observations. It distinguishes observed input data, transferred literature assumptions and generated simulation events. The implemented study links finite arrival demand to persistent aircraft identities, exports object-centric event logs and compares service-time, policy and capacity scenarios. The results reported here describe the completed simulation pilot as of 13 November 2026; they do not constitute a validated digital twin of BER.

\section{Background}
Airport Collaborative Decision Making (A-CDM) supports the exchange of timely operational information between airport partners and the Network Manager, particularly around turnaround and pre-departure processes \cite{eurocontrol_acdm}. However, access to an A-CDM-related application does not imply access to every internal ground-service event. A boundary interval can contain several concurrent activities, waiting and parking. Its length alone does not reveal the duration of cleaning, fueling or boarding.

Object-centric event logs retain connections between events and multiple objects rather than imposing one case identifier on every event. OCEL 2.0 also supports qualified relationships and changes in object attributes \cite{ocel2}. These connections are useful for airport operations: one turnaround involves an aircraft and a stand, while reusable service resources may participate in many turnarounds.

The executable model is derived from the \emph{Airport Ground Handling (Extended)} example provided by the supervisor, Istv\'an Koren, in OCPN Studio \cite{ocpn_studio}. The inherited process topology and resource logic are explicitly attributed to that model. The student's contribution is the empirical demand adapter, traceable pairing and assumption management, literature-informed timing changes, event-log analysis and controlled comparison of scenarios.

\section{Motivation}
The motivation is to explain simulated waiting and readiness through interacting activities and resources. Adding capacity may remove a local queue without substantially changing the aircraft's final completion time. Conversely, a policy rule can create a long apparent service delay without a resource shortage. Distinguishing these mechanisms is necessary before drawing conclusions from a turnaround-time distribution.

The study also examines how much can be learned from incomplete operational observations. Available data support arrival demand and candidate aircraft characteristics, but do not identify all internal service durations. A useful model must expose this distinction rather than treat generated service events as additional empirical evidence.

\section{Problem statement}
The original research objective is to use A-CDM information to calibrate an object-centric turnaround simulation and investigate resource allocation. In the available collection, internal service start/end times and measured block times are absent. The implemented study therefore combines OpenSky boundary proxies with A-CDM reference records, and treats transferred service timings and unobserved operational context as explicit assumptions. Calibration is currently limited to evidence-supported input construction; no fitted predictive service-time model or completed empirical validation is claimed.

\section{Research questions}
The original research questions are retained:
\begin{description}
\item[RQ1:] How can A-CDM Data be used to calibrate activity durations and temporal dependencies in an OCPN Simulation?
\item[RQ2:] How do dependencies between aircraft, stands, and shared ground-service resources contribute to waiting times and delay propagation at terminal stands and remote stands in an OCPN simulation?
\item[RQ3:] How can the calibrated simulation be used to evaluate resource-allocation strategies and identify configurations that reduce turnaround delays and resource contention?
\end{description}

RQ1 is addressed by an observability audit and evidence-informed model inputs, with clear limits on activity-level calibration. RQ2 is investigated through shared-resource waiting, parallel-branch readiness and object-centric extraction. A controlled terminal-versus-remote stand comparison and a quantified inter-flight propagation analysis remain open. RQ3 is addressed by repeated conditional capacity experiments, without claiming optimal real BER staffing. Predictive accuracy evaluation is a methodological requirement across these questions, rather than a replacement definition of RQ2.

\section{Contributions and current scope}
The completed work provides a reproducible preparation pipeline for aircraft-linked interval candidates; daily executable OCPN derivatives; OCEL exports inspected in Promenade; and a 60-replication pilot separating timing assumptions, inherited policy and resource capacity. It also identifies important threats to validity: boundary proxies differ from block times, long stays require departure-scheduling logic, and inherited policy is evaluated on the wrong local-time basis. Chapters~\ref{ch:methods}--\ref{ch:discussion} document both the completed results and the work still required before operational conclusions.
